\documentclass[titlepage,a4paper]{jsarticle}
\usepackage{../../../sty/import}% 各種パッケージインポート
\usepackage{../../../sty/title}% タイトルページの変更

%% タイトルページの変数
% レポートタイトル
\title{心理学特論第四回目出席票}
% 提出日
\expdate{\today}
% 科目名
\subject{}
% 分野
\class{情報経営システム工学分野}
% 学年
\grade{M1}
% 学籍番号
\mynumber{24336488}
% 記述者
\author{本間三暉}

\begin{document}
% titleページ作成
\maketitle
\section*{keyword\#1}
ダニング・クルーガー効果
\section*{keyword\#2}
自己中心性バイアス

\section{心理学特論の第04回テーマの「ポジティブ・イリュージョン」，ざっくりいうと認知の歪みですが，わかっていてもなかなか是正できないものでもあります．以下の質問に答えてください．}
\subsection{私が正当に評価されていない！と評価した局面は？}
先週に引き続いて，友人と家でVCを繋いでゲームをしていた時の話だが，明らかに自分の方が正しいことを言っているのに，間違ってるのはお前だと言われたとき．
今でもそのゲーム内の判断が間違っていたとは思っていないが，自分の提案を頑なに受け入れられないことに対して，正当に評価されていないと感じた．

\subsection{本日の資料の中の「PDスタディ」の台詞について，あなたならどのように答えますか？理不尽な物言いに対してあたなが返す台詞とその理由を完結にわかりやすく書いてください．}
ある程度は反論するが，埒が明かないと思ったら口では謝りつつ，それ以降は相手の話を聞かないようにして自らの怒りを鎮める．
基本的にはそれで話が終わりなのでそれ以降その話題は出さないで過ごす．

\section{心理学特論課題サイコロワーク1 260430 回答記入欄}
\begin{enumerate}
  %% オリエンテーション見てね〜〜
  % (01) やる気や関心の深さ，のめりこみの度合いなどを示すもの．【カタカナ7文字で】
  \item コミットメント
  % (02) 楽しい気分だと楽しい記憶が，悲しい気分だと悲しい記憶を思い出す．【漢字4文字で】
  \item 文脈効果
  % (03) 目の前にいない時にその人のことを考えると，その人に対する感情が強化される．＜ヒント＞ ○○化と言います．【漢字3文字で】
  \item 結晶化
  % (04) 失敗経験によって，努力しても事態を好転させられないと思いこんでしまった状態．＜ヒント＞ ○○○無気力と言います．【漢字6文字で】
  \item 学習性無気力
  % (05) 他人との関係で，つくしたりつくされたりするのは公平に近いのがベストという考え．＜ヒント＞ ○○理論と言います．【漢字4文字で】
  \item 公平理論
  % (06) 人は得する場合は確実な方を，損する場合は不確実な方を選択するという考え．＜ヒント＞ ○○○○○○理論と言います．【8文字で】
  \item プロスペクト理論
  % (07) 関係ないことまで自分に関係あるように思え，他人の目や口が気になること．【8文字で】
  \item 自己標的バイアス
  % (08) 成功がもたらすストレスを避けるため，わざと成功しないような行動に出てしまう．＜ヒント＞ ○○○○動機と言います．【漢字6文字で】
  \item 成功回避動機
  % (09) 自己評価の高い人は自分を高く評価してくれる人を好み，低い人はその逆となること．＜ヒント＞ ○○○○○○理論と言います．【漢字8文字で】
  \item 認知的斉合性理論
  % (10) プラスの情報よりもマイナスの情報の方が人の記憶に残りやすく大きな効果を持つ．＜ヒント＞ ○効果と言います．【漢字3文字で】
  \item 負効果
  %% 
  %% 該当する適切な言葉をえらんでね〜
  % 人はみんな自己愛が強く，自分は特別です．そのため自分を過大評価する傾向があります．
  % 例えば，自分も頑張ったが，同僚も同じように頑張った状況で，両者の評価に差がない場合に，
  %「なぜあんなに頑張ったのにそれが評価されないのだ」と不満が湧きます．相手も同じです．
  % このように自分を過大評価する心理傾向のことを（11）といいます．心理学者（12）たちの研究では，
  % 運動能力に対して60％の人が自分は平均より上とみなしており，自分は平均より下とみなす人は6％しかいませんでした．
  % 平均を上回る人が60％もいるのは統計的にはあり得ないことで，運動能力のように客観的にわかりやすい性質に関してもここまで（13）が歪んでいるので，
  %（14）やリーダーシップでは，自己能力がさらに自分に都合よく歪んでいると推測されます．リーダーシップでは，70％が自分は平均より上とみなし，
  % 平均より下とみなす人は2％しかいませんでした．基準がはっきりしないため，
  % 実際はうまくやれてなくても自分は上手くやっているといった思い込みを抱きやすいと考えられます．別の調査データでも，管理職の90％が，
  % 自分の能力は他の管理職より優れているとみなしていることや，大学教授の94％が自分を平均より優れた業績を上げているとみなしていることが示されています．
  % こうしてみるといかに多くの人が自分を過大評価しているかわかります．たとえ正当な評価が行われたとしても多くの人が，自分は正当に評価されていないと感じてしまうのです．
  % (15) のシステムを整えれば不満がなくなるというわけではないのです．
  % ＝ (11)～(15) 解答群 ＝
  % ＊問題に直接関係のない言葉も含まれています．
  % コミュニケーション
  % ネガティブイリュージョン
  % ポジティブイリュージョン
  % バイアス
  % 自信過剰
  % 自己認知
  % 人事評価
  % 他己認知
  % ダニング
  % ワトソン
  \item ポジティブイリュージョン
  \item ダニング
  \item 自己認知
  \item コミュニケーション
  \item 人事評価
  % 学習心理学の用語では，いわゆる「やる気」は動機付けの問題として研究されてきた．学ぶこと自体が面白く興味深いのは，(16) 動機付けによる学習行動である．
  % 教師が指導方法を工夫するだけでなく，直接学生をほめたり叱ったりすることも動機付けを高める．
  % これは(17) 動機付けであり，時と場合に応じた動機付けを適切に組み合わせることが大切である．

  % A.バンデューラはやる気と自信について，(18) という用語を使って説明している．
  % (18) とは，ある結果を得るために自分が発揮できると考えたパワーのことである．
  % ある課題を達成しようとするときの，自分の (19) に対する見積もり，自信，信念である．(18) が高い学生は成し遂げようと行動するパワーが強く，失敗しても挫折しにくい．
  % 他方で (18) の低い学生は課題に不安を持ち，すぐに諦める．従って，学生の (18) を高めることが，やる気を引き出すポイントになる．
  % A.バンデューラは (18) を育成するために4つのポイントを挙げている．それらは，行動の達成・代理的体験・言語的説得・(20) である．

  % ＝ (16)～(20) 解答群 ＝
  % ＊問題に直接関係のない言葉も含まれています．
  % 欲望の喚起
  % 情動の喚起
  % 効力
  % 能力
  % 内発的
  % 外発的
  % 代理的
  % 言語的
  % 自己効力感
  % 自己達成感
  \item 内発的
  \item 外発的
  \item 自己効力感
  \item 能力
  \item 情動の喚起
\end{enumerate}
% 参考文献
\begin{thebibliography}{99}
\bibitem{} 心理学特論 配布資料
\bibitem{} 心理学特論 授業スライド
\end{thebibliography}

\end{document}